\begin{proof}[Proof of \cref{obj:lem:calibrated-support}]
We use the calibration maps and families of
\cref{obj:def:calibration-maps,obj:def:continuous-calibrated-priors}.
All finite coordinate spaces below carry the maximum norm, and total derivatives carry the induced multilinear operator norm.

\begin{enumerate}
\item \textit{Local formulas and smoothness radii.}
By \cref{obj:lem:exact-calibrations}, choose a calibration radius
\(\varepsilon_{\mathrm{cal}}>0\) on which the stated smoothness, endpoint, and matching identities hold.

For a local sign pair \(\mathbf s=(s_1,s_2)\in\{-1,1\}^2\), write
\[
Z_u(\mathbf s)
=s_1\cos(\pi u/2)+s_2\sin(\pi u/2).
\]
For \(a,y\in\{0,1\}\), define the local mixed cells by
\[
\begin{gathered}
e^{\mathrm{loc}}_{b,\mathbf s}(\eta,\zeta,u)
=\frac12+(1-2b)\eta\mathfrak q(\eta,\zeta,u)Z_u(\mathbf s),
\qquad b\in\{0,1\},\\
m^{\mathrm{loc}}_{\mathbf s}(\zeta,u)
=\frac12+\zeta Z_u(\mathbf s),\\
c^{\mathrm{loc}}_{0,\mathbf s}(\eta,\zeta,u)=0,\qquad
c^{\mathrm{loc}}_{1,\mathbf s}(\eta,\zeta,u)
=\mathfrak c\!\left(
32\eta\zeta,e^{\mathrm{loc}}_{1,\mathbf s}(\eta,\zeta,u),
m^{\mathrm{loc}}_{\mathbf s}(\zeta,u)\right),
\end{gathered}
\]
and
\[
\begin{aligned}
\pi^{\mathrm{mix}}_{b,\mathbf s,11}
&=e^{\mathrm{loc}}_{b,\mathbf s}m^{\mathrm{loc}}_{\mathbf s}
  +c^{\mathrm{loc}}_{b,\mathbf s},&
\pi^{\mathrm{mix}}_{b,\mathbf s,10}
&=e^{\mathrm{loc}}_{b,\mathbf s}(1-m^{\mathrm{loc}}_{\mathbf s})
  -c^{\mathrm{loc}}_{b,\mathbf s},\\
\pi^{\mathrm{mix}}_{b,\mathbf s,01}
&=(1-e^{\mathrm{loc}}_{b,\mathbf s})m^{\mathrm{loc}}_{\mathbf s}
  -c^{\mathrm{loc}}_{b,\mathbf s},&
\pi^{\mathrm{mix}}_{b,\mathbf s,00}
&=(1-e^{\mathrm{loc}}_{b,\mathbf s})(1-m^{\mathrm{loc}}_{\mathbf s})
  +c^{\mathrm{loc}}_{b,\mathbf s}.
\end{aligned}
\]
The arguments of these cells are \((\eta,\zeta,u)\).

Use \(\mathrm r\) and \(\mathrm c\) for the random and comparator fair branches, respectively. Their local risks and cells are
\[
\begin{aligned}
\mu^{\mathrm{loc}}_{0,\mathrm r}(t,\delta,u,\mathbf s)
&=\mathfrak p(t,\delta,u)+\delta Z_u(\mathbf s),&
\mu^{\mathrm{loc}}_{1,\mathrm r}
&=\mathsf G(t,\mu^{\mathrm{loc}}_{0,\mathrm r}),\\
\mu^{\mathrm{loc}}_{0,\mathrm c}(t,\delta,u,\mathbf s)
&=\mathfrak p(t,\delta,u),&
\mu^{\mathrm{loc}}_{1,\mathrm c}
&=\mathsf G(\mathsf T(t,\delta),\mu^{\mathrm{loc}}_{0,\mathrm c}),\\
\pi^{\mathrm{fair}}_{\iota,\mathbf s,ay}(t,\delta,u)
&=\frac12
\begin{cases}
\mu^{\mathrm{loc}}_{a,\iota}(t,\delta,u,\mathbf s),&y=1,\\
1-\mu^{\mathrm{loc}}_{a,\iota}(t,\delta,u,\mathbf s),&y=0,
\end{cases}
&&\iota\in\{\mathrm r,\mathrm c\}.
\end{aligned}
\]
The comparator expressions are independent of \(\mathbf s\).

The smooth-neighborhood conclusion of
\cref{obj:lem:exact-calibrations} supplies positive radii
\(r_{\mathrm{mix},\infty}\) and \(r_{\mathrm{fair},\infty}\) on whose closed parameter regions all these local formulas are smooth at every point. Set
\[
\rho=\frac12\min\{r_{\mathrm{mix},\infty},
                         r_{\mathrm{fair},\infty}\}>0.
\]
Thus, whenever \(0<\varepsilon\le\rho\), the local mixed formulas are smooth at every
\[
|\eta|<2\varepsilon,\qquad |\zeta|<2\varepsilon,\qquad u\in[0,1],
\]
and the local fair formulas are smooth at every
\[
t\in[0,1/4],\qquad |\delta|<2\varepsilon,\qquad u\in[0,1].
\]

For preliminary estimates, interpret a supplied spatial cell profile
\(\varpi=(\varpi_{ay})_{a,y\in\{0,1\}}\) through the convention
\[
\varpi^\dagger_{ay}(x)=
\begin{cases}
\varpi_{ay}(x),
&\text{if all cells are continuous, }
  0<\varpi_{ay}(x)<1,\ \sum_{a,y}\varpi_{ay}(x)=1
  \text{ for every }x,\\
1/4,&\text{otherwise}.
\end{cases}
\]
In either case the covariate law is \(\lambda\).
We will establish the first case throughout the required neighborhoods.
% lean: CausalSmith.Stat.LogoddsLowsmoothFrontier.calibrated_density_smoothness_of_valid

\item \textit{A uniform bound for amplitude derivatives.}
The local derivative bounds supplied by
\cref{obj:lem:exact-calibrations} give positive constants
\(B_{\mathrm{mix},1},B_{\mathrm{mix},2},
 B_{\mathrm{fair},1},B_{\mathrm{fair},2}\)
and positive radii \(r_{\mathrm{mix},D},r_{\mathrm{fair},D}\) such that, for \(m=1,2\),
\[
\begin{aligned}
\left\|D^m\pi^{\mathrm{mix}}_{b,\mathbf s,ay}
                   (\eta,\zeta,u)\right\|
&\le B_{\mathrm{mix},m},
&&|\eta|,|\zeta|\le r_{\mathrm{mix},D},\quad u\in[0,1],\\
\left\|D^m\pi^{\mathrm{fair}}_{\iota,\mathbf s,ay}
                   (t,\delta,u)\right\|
&\le B_{\mathrm{fair},m},
&&t\in[0,1/4],\quad |\delta|\le r_{\mathrm{fair},D},
  \quad u\in[0,1].
\end{aligned}
\]
These are total derivatives in the three displayed arguments, uniformly over the finite indices. Define
\[
r_D=\min\{r_{\mathrm{mix},D},r_{\mathrm{fair},D}\},
\qquad
M=4\bigl(B_{\mathrm{mix},1}+B_{\mathrm{mix},2}
        +B_{\mathrm{fair},1}+B_{\mathrm{fair},2}\bigr)+1.
\]
In particular, \(r_D>0\) and \(M\ge1\).

The amplitude slices have linear parts
\[
L_{\mathrm{mix}}(\eta,\zeta)=(\eta,\zeta,0),
\qquad
L_{\mathrm{fair}}(\delta)=(0,\delta,0),
\qquad
\|L_{\mathrm{mix}}\|=\|L_{\mathrm{fair}}\|=1.
\]
For a smooth function, the \(m\)th derivative of its composition with an affine map is its \(m\)th derivative evaluated on the linear part in each argument. Consequently, fixing \(u\), or fixing \(t,u\), gives
\[
\begin{aligned}
\left\|D^m_{(\eta,\zeta)}
      \bigl[4\pi^{\mathrm{mix}}_{b,\mathbf s,ay}(\eta,\zeta,u)\bigr]\right\|
&\le4B_{\mathrm{mix},m}\le M,\\
\left\|D^m_\delta
      \bigl[4\pi^{\mathrm{fair}}_{\iota,\mathbf s,ay}(t,\delta,u)\bigr]\right\|
&\le4B_{\mathrm{fair},m}\le M,
\qquad m=1,2.
\end{aligned}
\]
The factor four comes from density relative to the uniform measure on the four labels. These bounds involve the fully substituted formulas, including the comparator effect.
% lean: CausalSmith.Stat.LogoddsLowsmoothFrontier.calibrated_local_density_derivative_envelopes

\item \textit{Fair prognosis estimates.}
By \cref{obj:lem:exact-calibrations}, choose
\(\varepsilon_{\mathrm f}>0\) and \(B_{\mathrm p}>0\) such that
\[
\left|\mathfrak p(t,\delta,u)-\frac25\right|
+\left|\partial_u\mathfrak p(t,\delta,u)\right|
\le B_{\mathrm p}\delta^2
\]
for \(t\in[0,1/4]\), \(|\delta|\le\varepsilon_{\mathrm f}\), and \(u\in[0,1]\).
Put
\[
r_E=r_O
=\min\left\{\varepsilon_{\mathrm f},\frac1{100},
                    \frac1{100(B_{\mathrm p}+1)}\right\},
\qquad
r_L=\min\{r_E,r_{\mathrm{fair},\infty}\}.
\]
All three radii are positive, and
\[
B_{\mathrm p}r_E\le\frac1{100}.
\]
Hence, for \(|\delta|\le r_E\),
\[
\left|\mathfrak p(t,\delta,u)-\frac25\right|
\le B_{\mathrm p}\delta^2
\le B_{\mathrm p}r_E^2
\le\frac1{100}.
\]
Since \(|Z_u(\mathbf s)|\le2\), both fair control risks lie in
\([1/3,2/3]\): the random risk differs from \(2/5\) by at most
\(1/100+2/100\), and the comparator risk by at most \(1/100\).

We record the logit estimates used here and below. For
\(\xi_1,\xi_2\in[1/3,2/3]\), the mean value theorem applied separately to the logarithms of the risks and their complements gives
\[
|\log \xi_1-\log \xi_2|\le3|\xi_1-\xi_2|,
\qquad
|\log(1-\xi_1)-\log(1-\xi_2)|\le3|\xi_1-\xi_2|.
\]
Thus
\[
|\operatorname{logit}(\xi_1)-\operatorname{logit}(\xi_2)|
\le6|\xi_1-\xi_2|.
\]
Also, with
\[
\omega_\xi=\frac{\xi}{1-\xi},
\qquad \xi\in[1/3,2/3],
\]
we have \(\omega_\xi\in[1/2,2]\). The elementary inequalities
\[
1-\omega_\xi^{-1}\le\log\omega_\xi\le\omega_\xi-1
\]
therefore imply
\[
|\operatorname{logit}(\xi)|\le1.
\]
This proves the fair prognosis envelope on \(|\delta|\le r_E\).

For a valid fair profile, comparison with the fixed risk \(2/5\) yields
\[
\left|\nu(P,x)-\operatorname{logit}\!\left(\frac25\right)\right|
\le6(B_{\mathrm p}\delta^2+2|\delta|).
\]
The triangle inequality consequently gives
\[
|\nu(P,x)-\nu(P,z)|
\le12B_{\mathrm p}\delta^2+24|\delta|.
\]
For an invalid preliminary profile, its replacement has prognosis logit zero, so the same oscillation inequality holds.

For any auxiliary radius \(0<\varepsilon_{\mathrm{amp}}\le r_O\), any
\(\beta\ge0\), and any \(k\ge1\), set
\[
\delta_{\mathrm{amp}}
=\varepsilon_{\mathrm{amp}}k^{-\beta}.
\]
Because \(0<k^{-\beta}\le1\),
\[
0\le\delta_{\mathrm{amp}}\le\varepsilon_{\mathrm{amp}}\le r_O,
\qquad
\delta_{\mathrm{amp}}^2\le r_O|\delta_{\mathrm{amp}}|.
\]
It follows that
\[
\begin{aligned}
12B_{\mathrm p}\delta_{\mathrm{amp}}^2
   +24|\delta_{\mathrm{amp}}|
&\le(12B_{\mathrm p}r_O+24)|\delta_{\mathrm{amp}}|\\
&\le26|\delta_{\mathrm{amp}}|
\le2k^{-\beta},
\end{aligned}
\]
where the last inequality uses \(26\varepsilon_{\mathrm{amp}}\le2\).

For the linear estimate, define the within-cell coordinate by
\[
j_k(x)=\min\{\lfloor kx\rfloor,k-1\},
\qquad
u_k(x)=kx-j_k(x),
\qquad x\in[0,1].
\]
Here \(0\le j_k(x)<k\) and \(u_k(x)\in[0,1]\): below the last cell this follows from the floor inequalities, and in the last cell it follows from
\(k-1\le kx\le k\). In particular this convention includes \(x=1\) and \(k=1\).

Within a cell, the elementary sine and cosine difference bounds give
\[
|Z_u(\mathbf s)-Z_v(\mathbf s)|
\le\pi|u-v|\le4|u-v|,
\qquad u,v\in[0,1].
\]
Adjacent profiles agree at the shared sign. Subdividing a segment at its cell boundaries therefore gives
\[
|\mathsf Z_{k,\sigma}(x)-\mathsf Z_{k,\sigma}(z)|
\le4k|x-z|.
\]
Similarly, smoothness of the fair root, its derivative bound, and
\(\mathfrak p(t,\delta,0)=\mathfrak p(t,\delta,1)=2/5\) give
\[
|\mathfrak p(t,\delta,u_k(x))
 -\mathfrak p(t,\delta,u_k(z))|
\le B_{\mathrm p}\delta^2k|x-z|.
\]
Thus each fair control risk satisfies
\[
|\mu_0(x)-\mu_0(z)|
\le(B_{\mathrm p}\delta^2+4|\delta|)k|x-z|.
\]
Applying the logit difference bound, or using the zero logit in the replacement case, proves
\[
|\nu(P,x)-\nu(P,z)|
\le6(B_{\mathrm p}\delta^2+4|\delta|)k|x-z|.
\]
For \(0<\varepsilon_{\mathrm{amp}}\le r_L\),
\[
6(B_{\mathrm p}\delta_{\mathrm{amp}}^2
   +4|\delta_{\mathrm{amp}}|)
\le25|\delta_{\mathrm{amp}}|,
\qquad
25\varepsilon_{\mathrm{amp}}\le1.
\]
Since \(k^{1-\beta}=k^{-\beta}k\), we obtain
\[
|\nu(P,x)-\nu(P,z)|
\le k^{1-\beta}|x-z|.
\]
These estimates apply to both fair branches uniformly over
\(t\in[0,1/4]\).
% lean: CausalSmith.Stat.LogoddsLowsmoothFrontier.fairCells_uniform_prognosis_envelope
% lean: CausalSmith.Stat.LogoddsLowsmoothFrontier.fairCells_uniform_prognosis_oscillation
% lean: CausalSmith.Stat.LogoddsLowsmoothFrontier.fairCells_uniform_prognosis_linear

\item \textit{Mixed propensity and alternative-prognosis estimates.}
The mixed-root smoothness and derivative bounds in
\cref{obj:lem:exact-calibrations} supply a positive radius \(r_{\mathrm q}\) and a constant \(B_{\mathrm q}>0\) such that
\[
|\partial_u\mathfrak q(\eta,\zeta,u)|\le B_{\mathrm q}
\]
for \(|\eta|,|\zeta|\le r_{\mathrm q}\), \(u\in[0,1]\).
Indeed, the spatial derivative is the total first derivative evaluated on the unit vector in the spatial coordinate.

The endpoint identity for \(\mathfrak q\), followed by the mean value theorem within each cell and subdivision across cells, gives
\[
|\mathfrak q(\eta,\zeta,u_k(x))
 -\mathfrak q(\eta,\zeta,u_k(z))|
\le B_{\mathrm q}k|x-z|.
\]
Together with \(|\mathfrak q|\le5/4\), \(|\mathsf Z_{k,\sigma}|\le2\), and the field estimate above, the product decomposition gives
\[
\begin{aligned}
&|\mathfrak q(\eta,\zeta,u_k(x))\mathsf Z_{k,\sigma}(x)
 -\mathfrak q(\eta,\zeta,u_k(z))\mathsf Z_{k,\sigma}(z)|\\
&\quad\le
\frac54|\mathsf Z_{k,\sigma}(x)-\mathsf Z_{k,\sigma}(z)|
+2|\mathfrak q(\eta,\zeta,u_k(x))
      -\mathfrak q(\eta,\zeta,u_k(z))|\\
&\quad\le(5+2B_{\mathrm q})k|x-z|.
\end{aligned}
\]
For \(|\eta|,|\zeta|\le1/100\), the mixed margins satisfy
\[
|e_b(x)-1/2|\le\frac52|\eta|\le\frac1{40},
\qquad
|m_\sigma(x)-1/2|\le2|\zeta|\le\frac1{40}.
\]
Their propensity margins consequently belong to \([1/3,2/3]\). For a valid profile, summing the treated cells identifies its propensity with \(e_b\), and hence
\[
|g(P_{b,\sigma},x)-g(P_{b,\sigma},z)|
\le6|\eta|(5+2B_{\mathrm q})k|x-z|.
\]
The replacement profile has propensity logit zero, giving the same bound. Define
\[
r_P=\min\left\{r_{\mathrm q},\frac1{100},
                  \frac1{6(5+2B_{\mathrm q})}\right\}>0.
\]
For
\[
\eta_{\mathrm{amp}}
=\varepsilon_{\mathrm{amp}}k^{-\alpha},
\qquad
\zeta_{\mathrm{amp}}
=\varepsilon_{\mathrm{amp}}k^{-\beta},
\]
with \(0<\varepsilon_{\mathrm{amp}}\le r_P\) and
\(\alpha,\beta\ge0\), we therefore have
\[
|g(P_{b,\sigma},x)-g(P_{b,\sigma},z)|
\le k^{1-\alpha}|x-z|.
\]

For the alternative prognosis, consider the control-logit map
\[
\mathcal F_{\mathrm{ctrl}}(\vartheta,\xi,\upsilon)
=\operatorname{logit}\!\left(
\upsilon-\frac{\mathfrak c(\vartheta,\xi,\upsilon)}{1-\xi}
\right).
\]
It is smooth near \((0,1/2,1/2)\): the covariance branch is smooth there, the denominator equals \(1/2\) at the center, and the argument of the logit equals \(1/2\). A \(C^1\) function is locally Lipschitz. Choose a neighborhood
\(\mathcal U_{\mathrm{ctrl}}\), a Lipschitz constant
\(K_{\mathrm{ctrl}}\ge0\), and \(h_{\mathrm{ctrl}}>0\) such that
\[
\left\{(\vartheta,\xi,\upsilon):
\max\{|\vartheta|,|\xi-1/2|,|\upsilon-1/2|\}
<h_{\mathrm{ctrl}}\right\}
\subseteq\mathcal U_{\mathrm{ctrl}}.
\]
Set
\[
r_{\mathrm{ctrl}}=h_{\mathrm{ctrl}}/2,
\qquad C_{\mathrm{ctrl}}=K_{\mathrm{ctrl}}+1.
\]
Comparing points with the same first coordinate gives
\[
\begin{aligned}
&|\mathcal F_{\mathrm{ctrl}}(\vartheta,\xi_1,\upsilon_1)
 -\mathcal F_{\mathrm{ctrl}}(\vartheta,\xi_2,\upsilon_2)|\\
&\qquad\le C_{\mathrm{ctrl}}
       (|\xi_1-\xi_2|+|\upsilon_1-\upsilon_2|)
\end{aligned}
\]
whenever
\[
|\vartheta|,\ |\xi_1-1/2|,\ |\xi_2-1/2|,
\ |\upsilon_1-1/2|,\ |\upsilon_2-1/2|
\le r_{\mathrm{ctrl}}.
\]

Define
\[
r_N=\min\left\{r_{\mathrm q},\frac1{100},
 \frac{r_{\mathrm{ctrl}}}{3},
 \frac1{C_{\mathrm{ctrl}}(9+2B_{\mathrm q})}\right\}>0.
\]
For \((\alpha,\beta)\in\mathcal D\), we have
\(0<\beta<\alpha\le1\). At the rate-scaled amplitudes with
\(0<\varepsilon_{\mathrm{amp}}\le r_N\),
\[
0\le\eta_{\mathrm{amp}}\le\zeta_{\mathrm{amp}}
\le r_N.
\]
Moreover,
\[
32|\eta_{\mathrm{amp}}\zeta_{\mathrm{amp}}|
\le\frac{32}{100}r_N\le r_{\mathrm{ctrl}},
\qquad
\frac52|\eta_{\mathrm{amp}}|\le r_{\mathrm{ctrl}},
\qquad
2|\zeta_{\mathrm{amp}}|\le r_{\mathrm{ctrl}}.
\]
The control-logit comparison therefore applies to the alternative margins. The two margin differences obey
\[
\begin{aligned}
|e_1(x)-e_1(z)|
&\le|\eta_{\mathrm{amp}}|(5+2B_{\mathrm q})k|x-z|,\\
|m_\sigma(x)-m_\sigma(z)|
&\le4|\zeta_{\mathrm{amp}}|k|x-z|.
\end{aligned}
\]
For a valid table, its control risk is
\[
\mu_0(x)=m_\sigma(x)-\frac{c_1(x)}{1-e_1(x)}.
\]
Using \(\eta_{\mathrm{amp}}\le\zeta_{\mathrm{amp}}\) gives
\[
\begin{aligned}
|\nu(P_{1,\sigma},x)-\nu(P_{1,\sigma},z)|
&\le C_{\mathrm{ctrl}}(9+2B_{\mathrm q})
       \zeta_{\mathrm{amp}}k|x-z|\\
&\le k^{1-\beta}|x-z|.
\end{aligned}
\]
The replacement case again has zero prognosis logit.
% lean: CausalSmith.Stat.LogoddsLowsmoothFrontier.mixedCells_uniform_propensity_linear
% lean: CausalSmith.Stat.LogoddsLowsmoothFrontier.mixedCells_uniform_alternative_prognosis_linear

\item \textit{One radius and the remaining mixed spatial bounds.}
Choose
\[
R_{\mathrm{support}}
=\min\left\{\varepsilon_{\mathrm{cal}},\rho,r_D,r_E,r_O,r_L,
                         r_P,r_N,\frac1{100}\right\},
\qquad
\varepsilon=\frac12R_{\mathrm{support}}.
\]
Every entry in the minimum is positive. Therefore
\[
0<\varepsilon\le\frac1{100},
\qquad
2\varepsilon\le\varepsilon_{\mathrm{cal}},
\qquad
2\varepsilon\le\frac1{100},
\qquad
\varepsilon\le\rho,r_D,r_E,r_O,r_L,r_P,r_N.
\]

We supply the remaining mixed estimates at
\[
\eta=\varepsilon k^{-\alpha},
\qquad
\zeta=\varepsilon k^{-\beta}.
\]
For every nonnegative exponent \(\gamma\),
\[
0<k^{-\gamma}\le1,
\qquad
|\varepsilon k^{-\gamma}|=\varepsilon k^{-\gamma}\le\varepsilon.
\]
Thus \(|\eta|,|\zeta|\le1/100\).

For the null branch the control risk equals \(m_\sigma\) whenever the profile is valid. The logit and field bounds give
\[
|\nu(P_{0,\sigma},x)-\nu(P_{0,\sigma},z)|
\le24|\zeta|k|x-z|
\le k^{1-\beta}|x-z|,
\]
because \(24\varepsilon\le1\). In the replacement case the left side is zero.

For oscillation, comparison of the propensity margin with \(1/2\) gives
\[
|g(P_{b,\sigma},x)|
\le6|e_b(x)-1/2|\le15|\eta|,
\]
and hence
\[
|g(P_{b,\sigma},x)-g(P_{b,\sigma},z)|
\le30|\eta|
\le2k^{-\alpha}.
\]

We also need a covariance estimate for the mixed prognosis. For
\[
|\vartheta|,\ |\xi-1/2|,\ |\upsilon-1/2|\le\frac18,
\]
define
\[
\begin{gathered}
d_\vartheta=e^\vartheta-1,\qquad
v_{\xi,\upsilon}=\xi(1-\xi)\upsilon(1-\upsilon),\\
L_{\vartheta,\xi,\upsilon}
=1+d_\vartheta\{\xi(1-\upsilon)+(1-\xi)\upsilon\},\\
S_{\vartheta,\xi,\upsilon}
=\sqrt{L_{\vartheta,\xi,\upsilon}^2
              -4d_\vartheta^2v_{\xi,\upsilon}},\qquad
D_{\vartheta,\xi,\upsilon}
=L_{\vartheta,\xi,\upsilon}+S_{\vartheta,\xi,\upsilon}.
\end{gathered}
\]
The elementary exponential bound and the margin bounds give
\[
|d_\vartheta|\le2|\vartheta|\le\frac14,\qquad
0\le v_{\xi,\upsilon}\le\frac1{16},\qquad
0\le\xi(1-\upsilon)+(1-\xi)\upsilon\le1.
\]
Consequently,
\[
L_{\vartheta,\xi,\upsilon}\ge\frac34,\qquad
L_{\vartheta,\xi,\upsilon}^2
       -4d_\vartheta^2v_{\xi,\upsilon}\ge0,
\qquad
D_{\vartheta,\xi,\upsilon}\ge\frac34.
\]
The rationalized covariance formula now yields
\[
|\mathfrak c(\vartheta,\xi,\upsilon)|
D_{\vartheta,\xi,\upsilon}
=2|d_\vartheta|v_{\xi,\upsilon},
\]
so
\[
|\mathfrak c(\vartheta,\xi,\upsilon)|
\le\frac1{24},
\qquad
|\mathfrak c(\vartheta,\xi,\upsilon)|
\le|\vartheta|.
\]
The first inequality uses \(|d_\vartheta|\le1/4\); the second uses
\(|d_\vartheta|\le2|\vartheta|\).

For the mixed alternative,
\[
|\vartheta|=|32\eta\zeta|\le\frac18,
\]
and its margins lie within \(1/40\) of \(1/2\). Thus, for either branch,
\[
|c_b(x)|\le32|\eta||\zeta|,
\qquad |c_b(x)|\le\frac1{24}.
\]
Summing the control cells gives \(1-e_b\), so their risk satisfies
\[
\mu_0(x)=m_\sigma(x)-\frac{c_b(x)}{1-e_b(x)}.
\]
The bounds
\[
\frac{19}{40}\le e_b,m_\sigma\le\frac{21}{40},
\qquad |c_b|\le\frac1{24}
\]
imply
\[
\frac13(1-e_b)\le(1-e_b)m_\sigma-c_b
                 \le\frac23(1-e_b),
\]
by multiplication and the displayed rational bounds. Hence
\(\mu_0(x)\in[1/3,2/3]\). Since \(1-e_b\ge19/40>1/3\),
\[
|\mu_0(x)-1/2|
\le|m_\sigma(x)-1/2|+3|c_b(x)|
\le2|\zeta|+96|\eta||\zeta|.
\]
Applying the logit comparison with \(1/2\), and using \(|\eta|\le1/100\), gives
\[
|\nu(P_{b,\sigma},x)|
\le6(2|\zeta|+96|\eta||\zeta|)
\le18|\zeta|.
\]
Therefore
\[
|\nu(P_{b,\sigma},x)-\nu(P_{b,\sigma},z)|
\le36|\zeta|\le2k^{-\beta}.
\]
All these bounds also hold for the zero logits of a replacement profile.

Combining these estimates with the preceding linear estimates proves, for both mixed branches,
\[
\begin{aligned}
|g(P_{b,\sigma},x)-g(P_{b,\sigma},z)|
&\le k^{1-\alpha}|x-z|,&
|g(P_{b,\sigma},x)-g(P_{b,\sigma},z)|
&\le2k^{-\alpha},\\
|\nu(P_{b,\sigma},x)-\nu(P_{b,\sigma},z)|
&\le k^{1-\beta}|x-z|,&
|\nu(P_{b,\sigma},x)-\nu(P_{b,\sigma},z)|
&\le2k^{-\beta}.
\end{aligned}
\]
% lean: CausalSmith.Stat.LogoddsLowsmoothFrontier.mixedCells_uniform_null_prognosis_linear
% lean: CausalSmith.Stat.LogoddsLowsmoothFrontier.mixedCells_uniform_propensity_oscillation
% lean: CausalSmith.Stat.LogoddsLowsmoothFrontier.mixedCells_uniform_prognosis_oscillation

\item \textit{Continuity and validity on the doubled boxes.}
Fix \(k\ge1\) and \(\sigma\). On cell \(j\), use the local pair
\((\sigma_j,\sigma_{j+1})\) and coordinate \(u=kx-j\).
For
\[
|\eta|,|\zeta|\le2\varepsilon,
\qquad
t\in[0,1/4],\quad |\delta|\le2\varepsilon,
\]
the local cell formulas are continuous in \(u\), by
\cref{obj:lem:exact-calibrations} and
\(2\varepsilon\le\varepsilon_{\mathrm{cal}}\).

At a shared endpoint, the two sign fields take the same value:
\[
Z_1(\sigma_j,\sigma_{j+1})
=\sigma_{j+1}
=Z_0(\sigma_{j+1},\sigma_{j+2}).
\]
The endpoint identities
\[
\mathfrak q(\eta,\zeta,1)=\mathfrak q(\eta,\zeta,0),
\qquad
\mathfrak p(t,\delta,1)=\mathfrak p(t,\delta,0)=\frac25
\]
therefore make every mixed and fair cell agree across adjoining cells. Gluing the finitely many continuous cell profiles proves continuity on \([0,1]\), including its endpoints. When \(k=1\), there is one continuous profile and no interior boundary to glue.

Because \(2\varepsilon\le1/100\), the raw mixed cells satisfy
\[
\frac{19}{40}\le e_b,m_\sigma,1-e_b,1-m_\sigma\le\frac{21}{40},
\qquad |c_b|\le\frac1{24}.
\]
Each product entering a mixed cell lies in
\([361/1600,441/1600]\). Accordingly,
\[
4\left(\frac{361}{1600}-\frac1{24}\right)
\le4\pi^{\mathrm{mix}}_{b,\mathbf s,ay}
\le4\left(\frac{441}{1600}+\frac1{24}\right),
\]
which implies
\[
\frac12\le4\pi^{\mathrm{mix}}_{b,\mathbf s,ay}\le2.
\]

For the fair cells, the selected root belongs to \((3/10,1/2)\); this includes its default value \(2/5\). Since
\[
|\delta Z_u(\mathbf s)|\le\frac1{50},
\]
both fair control risks belong to \([1/4,3/5]\).
For \(t\in[0,1/4]\), the elementary exponential estimate gives
\[
0\le e^t-1\le\frac43t\le\frac13.
\]
The comparator effect lies between \(t/2\) and \(t\), by
\cref{obj:lem:exact-calibrations}. Thus the effect used in either fair branch lies in \([0,1/4]\). For
\(\vartheta\in[0,1/4]\) and \(\xi\in[1/4,3/5]\),
\[
\mathsf G(\vartheta,\xi)
=\frac{e^\vartheta\xi}{1+(e^\vartheta-1)\xi},
\qquad
\frac14\le\mathsf G(\vartheta,\xi)\le\frac34.
\]
To check the last bounds, the denominator is positive; multiplication by it reduces the lower bound to
\[
(1+e^\vartheta-1)\xi
\ge\frac14\{1+(e^\vartheta-1)\xi\},
\]
using \(\xi\ge1/4\), and the upper bound follows from
\(\xi\le3/5\) and \(e^\vartheta-1\le1/3\).
Consequently both fair arm risks and their complements are at least \(1/4\), and
\[
\frac12\le4\pi^{\mathrm{fair}}_{\iota,\mathbf s,ay}\le2.
\]

Finally, direct addition of either four-cell formula gives
\[
\sum_{a,y}\pi^{\mathrm{mix}}_{b,\mathbf s,ay}=1,
\qquad
\sum_{a,y}\pi^{\mathrm{fair}}_{\iota,\mathbf s,ay}=1.
\]
Every cell lies in \([1/8,1/2]\), so the cells are strictly interior. Together with continuity, this proves validity throughout both doubled closed boxes. The constructed laws therefore retain their supplied cells there.
% lean: CausalSmith.Stat.LogoddsLowsmoothFrontier.mixedCells_valid_of_continuous
% lean: CausalSmith.Stat.LogoddsLowsmoothFrontier.fairCells_valid_of_continuous

\item \textit{Transfer to actual densities and common open neighborhoods.}
Let \(\tau\) be the uniform probability measure on the four treatment--outcome labels:
\[
\tau(\{(a,y)\})=\frac14,\qquad a,y\in\{0,1\}.
\]
Define the actual one-record densities, relative to \(\lambda\otimes\tau\), by
\[
\begin{aligned}
\mathsf d^{\mathrm{mix}}_{b,k,\sigma,a,y,x}(\eta,\zeta)
&=4P_{b,\sigma}(A=a,Y=y\mid X=x),\\
\mathsf d^{\mathrm{fair}}_{\mathrm r,k,\sigma,t,a,y,x}(\delta)
&=4P_{t,\sigma}(A=a,Y=y\mid X=x),\\
\mathsf d^{\mathrm{fair}}_{\mathrm c,k,\sigma,t,a,y,x}(\delta)
&=4P_{*,t,\delta}(A=a,Y=y\mid X=x).
\end{aligned}
\]
The amplitudes of the laws on the right are those on the left.

Fix \(x\), and define its local sign pair by
\[
\mathbf s_{k,\sigma}(x)
=(\sigma_{j_k(x)},\sigma_{j_k(x)+1}).
\]
Validity on the doubled boxes identifies the actual densities there with the local slices:
\[
\begin{aligned}
\mathsf d^{\mathrm{mix}}_{b,k,\sigma,a,y,x}(\eta,\zeta)
&=4\pi^{\mathrm{mix}}_{b,\mathbf s_{k,\sigma}(x),ay}
                    (\eta,\zeta,u_k(x)),\\
\mathsf d^{\mathrm{fair}}_{\iota,k,\sigma,t,a,y,x}(\delta)
&=4\pi^{\mathrm{fair}}_{\iota,\mathbf s_{k,\sigma}(x),ay}
                    (t,\delta,u_k(x)).
\end{aligned}
\]
For the comparator the chosen pair is immaterial.

At a point of the smaller closed box, continuity of absolute value supplies a neighborhood lying inside the doubled box. Thus these identities hold on a neighborhood of every such point, and differentiation transfers the local bounds to the actual densities. Since \(\varepsilon\le r_D\),
\[
\begin{aligned}
\left\|D^m\mathsf d^{\mathrm{mix}}_{b,k,\sigma,a,y,x}
                            (\eta,\zeta)\right\|&\le M,
&&|\eta|,|\zeta|\le\varepsilon,\\
\left\|D^m\mathsf d^{\mathrm{fair}}_{\iota,k,\sigma,t,a,y,x}
                            (\delta)\right\|&\le M,
&&|\delta|\le\varepsilon,\quad t\in[0,1/4],
\end{aligned}
\qquad m=1,2.
\]

Moreover, set
\[
U_{\mathrm{mix}}
=\{(\eta,\zeta)\in\mathbb R^2:
                   |\eta|<2\varepsilon,\ |\zeta|<2\varepsilon\},
\qquad
U_{\mathrm{fair}}
=(-2\varepsilon,2\varepsilon).
\]
These sets are open and contain the required closed boxes. At every point of them, nearby validity identifies the actual density with a smooth local slice. Because \(\varepsilon\le\rho\), every mixed density is \(C^\infty\) on \(U_{\mathrm{mix}}\), and every fair density is \(C^\infty\) on \(U_{\mathrm{fair}}\), uniformly over the indicated indices and over \(t\in[0,1/4]\). In particular, they are \(C^2\) at every point of the smaller closed boxes.

For clarity, with
\[
\mathbf e_\eta=(1,0),\qquad \mathbf e_\zeta=(0,1),
\]
the total second-derivative bound includes
\[
\left|\partial_\eta\partial_\zeta
       \mathsf d^{\mathrm{mix}}(\eta,\zeta)\right|
=\left|D^2\mathsf d^{\mathrm{mix}}(\eta,\zeta)
                  [\mathbf e_\eta,\mathbf e_\zeta]\right|
\le M.
\]
It also applies to both first and second \(\delta\)-derivatives of the deterministic comparator.
% lean: CausalSmith.Stat.LogoddsLowsmoothFrontier.calibrated_density_derivative_envelopes_of_valid
% lean: CausalSmith.Stat.LogoddsLowsmoothFrontier.calibrated_density_smoothness_of_valid

\item \textit{Mixed model membership and effects.}
Fix \((\alpha,\beta)\in\mathcal D\), \(k\ge1\), and the rate-scaled amplitudes
\[
\eta=\varepsilon k^{-\alpha},\qquad
\zeta=\varepsilon k^{-\beta}.
\]
They lie in the valid smaller box. The constructed covariate marginal is uniform: for every Borel set \(\mathcal B_X\subseteq[0,1]\),
\[
P(X\in\mathcal B_X)
=\int_{\mathcal B_X}\sum_{a,y}p_{ay}(x)\,d\lambda(x)
=\lambda(\mathcal B_X).
\]

We identify the homogeneous effects from the positive cell odds. For the covariance formula introduced above, the identity
\[
S_{\vartheta,\xi,\upsilon}^2
=L_{\vartheta,\xi,\upsilon}^2
 -4d_\vartheta^2v_{\xi,\upsilon}
\]
and
\[
\mathfrak c(\vartheta,\xi,\upsilon)
=\frac{2d_\vartheta v_{\xi,\upsilon}}
       {D_{\vartheta,\xi,\upsilon}}
\]
give, by multiplication by the positive denominator,
\[
d_\vartheta\mathfrak c(\vartheta,\xi,\upsilon)^2
-L_{\vartheta,\xi,\upsilon}\mathfrak c(\vartheta,\xi,\upsilon)
+d_\vartheta v_{\xi,\upsilon}=0.
\]
For a table with margins \(\xi,\upsilon\) and covariance
\(\mathfrak c(\vartheta,\xi,\upsilon)\), expansion of its four cells then gives
\[
p_{11}p_{00}-e^\vartheta p_{10}p_{01}
=-d_\vartheta\mathfrak c(\vartheta,\xi,\upsilon)^2
 +L_{\vartheta,\xi,\upsilon}\mathfrak c(\vartheta,\xi,\upsilon)
 -d_\vartheta v_{\xi,\upsilon}=0.
\]
This computation includes \(\vartheta=0\). The null table, with covariance zero, directly satisfies
\(p_{11}p_{00}=p_{10}p_{01}\). Therefore the mixed branches satisfy
\[
\frac{p_{11}p_{00}}{p_{10}p_{01}}
=
\begin{cases}
1,&b=0,\\
e^{32\eta\zeta},&b=1.
\end{cases}
\]
All four cells are positive. Dividing by the positive treatment margins gives
\[
\mu_0=\frac{p_{01}}{p_{00}+p_{01}},
\qquad
\mu_1=\frac{p_{11}}{p_{10}+p_{11}},
\qquad
\frac{\mu_1/(1-\mu_1)}{\mu_0/(1-\mu_0)}
=\frac{p_{11}p_{00}}{p_{10}p_{01}}.
\]
Taking logarithms, and evaluating at covariate zero to identify the scalar effect, proves
\[
\theta(P_{0,\sigma})=0,\qquad
\theta(P_{1,\sigma})=32\eta\zeta,
\qquad
\mu_1(P_{b,\sigma},x)
=\ell\{\nu(P_{b,\sigma},x)+\theta(P_{b,\sigma})\}.
\]
The effect envelope follows from
\[
|32\eta\zeta|\le32\varepsilon^2\le\frac12.
\]
The propensity margin and the control risk lie in \([1/3,2/3]\), as established above, so
\[
\|g(P_{b,\sigma},\cdot)\|_\infty\le1,
\qquad
\|\nu(P_{b,\sigma},\cdot)\|_\infty\le1.
\]

To convert the spatial bounds into the required Hölder bounds, let
\(f:[0,1]\to\mathbb R\), \(0<\gamma\le1\), satisfy
\[
|f(x)-f(z)|\le k^{1-\gamma}|x-z|,
\qquad
|f(x)-f(z)|\le2k^{-\gamma}.
\]
Set
\[
d_{xz}=|x-z|\ge0.
\]
If \(kd_{xz}\le1\), then
\[
|f(x)-f(z)|
\le k^{-\gamma}(kd_{xz})
\le k^{-\gamma}(kd_{xz})^\gamma
=d_{xz}^\gamma.
\]
If \(kd_{xz}\ge1\), then
\[
|f(x)-f(z)|
\le2k^{-\gamma}
\le2k^{-\gamma}(kd_{xz})^\gamma
=2d_{xz}^\gamma.
\]
The first case includes \(x=z\), and both cases include \(\gamma=1\). Thus
\[
[f]_\gamma\le2.
\]
Apply this argument to \(g\) with \(\gamma=\alpha\) and to \(\nu\) with
\(\gamma=\beta\). Together with uniform design, homogeneity, and the three envelopes, these are precisely the conditions of
\cref{obj:def:model}. Hence both mixed laws belong to
\(\mathcal M(\alpha,\beta)\).
% lean: CausalSmith.Stat.LogoddsLowsmoothFrontier.mixedCells_model_bounds_of_spatial
% lean: CausalSmith.Stat.LogoddsLowsmoothFrontier.mixedCells_homogeneous_effect
% lean: CausalSmith.Stat.LogoddsLowsmoothFrontier.calibrated_model_of_bounds

\item \textit{Fair model membership and effects.}
Set
\[
\delta=\varepsilon k^{-\beta},
\qquad t\in[0,1/4].
\]
Both fair branches are valid. Their treatment probability is exactly \(1/2\), because the two cells in either treatment arm sum to \(1/2\). Therefore
\[
g(P_{t,\sigma},x)=g(P_{*,t,\delta},x)=0.
\]
This supplies the propensity envelope and both propensity spatial bounds.

The choices \(\varepsilon\le r_E,r_O,r_L\) and the fair estimates above give, for both branches,
\[
\|\nu(P,\cdot)\|_\infty\le1,\qquad
|\nu(P,x)-\nu(P,z)|\le k^{1-\beta}|x-z|,
\qquad
|\nu(P,x)-\nu(P,z)|\le2k^{-\beta}.
\]
The two-scale calculation in the preceding step gives
\[
[\nu(P,\cdot)]_\beta\le2.
\]

For any interior risk \(\xi\in(0,1)\), the positive denominator in the shift formula yields
\[
\mathsf G(\vartheta,\xi)(1-\xi)
=e^\vartheta\{1-\mathsf G(\vartheta,\xi)\}\xi.
\]
Applying this identity to the control risks of the two fair branches gives their positive cell cross-odds relations, and hence
\[
\theta(P_{t,\sigma})=t,\qquad
\theta(P_{*,t,\delta})=\mathsf T(t,\delta),
\qquad
\mu_1(P,x)=\ell\{\nu(P,x)+\theta(P)\}.
\]
By \cref{obj:lem:exact-calibrations},
\[
0\le\frac t2\le\mathsf T(t,\delta)\le t\le\frac14.
\]
Thus both effects satisfy the envelope \( |\theta(P)|\le1/2 \).
Uniform design follows by the same cell-normalization integral as above.
All conditions of \cref{obj:def:model} now hold for every random fair law and for its deterministic comparator.
% lean: CausalSmith.Stat.LogoddsLowsmoothFrontier.fairCells_model_bounds_of_prognosis
% lean: CausalSmith.Stat.LogoddsLowsmoothFrontier.fairCells_homogeneous_effect
% lean: CausalSmith.Stat.LogoddsLowsmoothFrontier.calibrated_model_of_bounds

\item \textit{Singleton identities at the rate-scaled amplitudes.}
Fix \(x\in[0,1]\), and put
\[
u=u_k(x),\qquad j=j_k(x).
\]
The indices \(j\) and \(j+1\) are distinct members of \(\{0,\ldots,k\}\), including when \(k=1\). For any function \(\Phi\) of their two signs, counting the \(2^{k-1}\) extensions of each pair gives
\[
2^{-(k+1)}\sum_{\sigma\in\{-1,1\}^{k+1}}
                 \Phi(\sigma_j,\sigma_{j+1})
=\frac14\sum_{s_1,s_2\in\{-1,1\}}\Phi(s_1,s_2).
\]
Consequently the full sign average of each cell is its local two-sign average.

Write \(E_{\mathbf s}\) for this uniform two-sign average:
\[
E_{\mathbf s}\Phi(\mathbf s)
=\frac14\sum_{\mathbf s\in\{-1,1\}^2}\Phi(\mathbf s).
\]
Symmetry and the trigonometric identity give
\[
E_{\mathbf s}Z_u=0,\qquad
E_{\mathbf s}Z_u^2
=\cos^2(\pi u/2)+\sin^2(\pi u/2)=1.
\]
The mixed calibration identity in
\cref{obj:lem:exact-calibrations} supplies
\[
E_{\mathbf s}\mathfrak c\!\left(
32\eta\zeta,\frac12-\eta\mathfrak q(\eta,\zeta,u)Z_u,
                 \frac12+\zeta Z_u\right)
=2\eta\zeta\mathfrak q(\eta,\zeta,u).
\]
Expanding the four local cells, using these moments and this identity, yields
\[
\begin{aligned}
E_{\mathbf s}\pi^{\mathrm{mix}}_{0,\mathbf s,11}
=E_{\mathbf s}\pi^{\mathrm{mix}}_{1,\mathbf s,11}
&=\frac14+\eta\zeta\mathfrak q(\eta,\zeta,u),\\
E_{\mathbf s}\pi^{\mathrm{mix}}_{0,\mathbf s,00}
=E_{\mathbf s}\pi^{\mathrm{mix}}_{1,\mathbf s,00}
&=\frac14+\eta\zeta\mathfrak q(\eta,\zeta,u),\\
E_{\mathbf s}\pi^{\mathrm{mix}}_{0,\mathbf s,10}
=E_{\mathbf s}\pi^{\mathrm{mix}}_{1,\mathbf s,10}
&=\frac14-\eta\zeta\mathfrak q(\eta,\zeta,u),\\
E_{\mathbf s}\pi^{\mathrm{mix}}_{0,\mathbf s,01}
=E_{\mathbf s}\pi^{\mathrm{mix}}_{1,\mathbf s,01}
&=\frac14-\eta\zeta\mathfrak q(\eta,\zeta,u).
\end{aligned}
\]
The cells of the actual laws equal these supplied cells by validity. Multiplying the equal full sign averages by \(2^{k+1}\) proves the asserted mixed sum identities.

For the fair family,
\cref{obj:lem:exact-calibrations} gives
\[
E_{\mathbf s}\mathsf G\!\left(
 t,\mathfrak p(t,\delta,u)+\delta Z_u\right)
=\mathsf G\!\left(\mathsf T(t,\delta),
                       \mathfrak p(t,\delta,u)\right).
\]
The control-risk average is \(\mathfrak p(t,\delta,u)\), by
\(E_{\mathbf s}Z_u=0\). Since treatment is fair, the four averaged cells are therefore
\[
\begin{aligned}
E_{\mathbf s}\pi^{\mathrm{fair}}_{\mathrm r,\mathbf s,01}
&=\frac12\mathfrak p(t,\delta,u),&
E_{\mathbf s}\pi^{\mathrm{fair}}_{\mathrm r,\mathbf s,00}
&=\frac12\{1-\mathfrak p(t,\delta,u)\},\\
E_{\mathbf s}\pi^{\mathrm{fair}}_{\mathrm r,\mathbf s,11}
&=\frac12\mathsf G(\mathsf T(t,\delta),\mathfrak p(t,\delta,u)),&
E_{\mathbf s}\pi^{\mathrm{fair}}_{\mathrm r,\mathbf s,10}
&=\frac12\{1-\mathsf G(\mathsf T(t,\delta),\mathfrak p(t,\delta,u))\}.
\end{aligned}
\]
These are exactly the comparator cells. Counting sign extensions and using validity proves the required full uniform-average identity. The calibration identities invoked here apply also at \(t=0\), at zero amplitudes, and at \(u=0,1\).
% lean: CausalSmith.Stat.LogoddsLowsmoothFrontier.mixedLaw_singleton_matching
% lean: CausalSmith.Stat.LogoddsLowsmoothFrontier.fairLaw_singleton_matching

\item \textit{Signed matching and completion.}
Now allow arbitrary signed amplitudes satisfying
\[
|\eta|\le\varepsilon,\qquad
|\zeta|\le\varepsilon,\qquad
|\delta|\le\varepsilon.
\]
They lie in the calibration region because
\(\varepsilon\le\varepsilon_{\mathrm{cal}}\), and their tables are valid by the doubled-box argument. The preceding local moment calculation and the matching identities of
\cref{obj:lem:exact-calibrations} apply on this whole signed region. Thus, for every \(a,y,x\),
\[
\sum_\sigma P_{0,\sigma}(A=a,Y=y\mid X=x)
=\sum_\sigma P_{1,\sigma}(A=a,Y=y\mid X=x),
\]
and, for every \(t\in[0,1/4]\),
\[
2^{-(k+1)}
\sum_\sigma P_{t,\sigma}(A=a,Y=y\mid X=x)
=P_{*,t,\delta}(A=a,Y=y\mid X=x).
\]

Finally, the raw cell bounds established on the doubled boxes transfer directly to the actual densities on the smaller boxes:
\[
\frac12\le
\mathsf d^{\mathrm{mix}}_{b,k,\sigma,a,y,x}(\eta,\zeta)
\le2,
\qquad
\frac12\le
\mathsf d^{\mathrm{fair}}_{\iota,k,\sigma,t,a,y,x}(\delta)
\le2.
\]
Together with the first- and second-derivative bounds, the common smooth open neighborhoods, and the rate-scaled membership and effect calculations, these prove every assertion. The radius \(\varepsilon\) and bound \(M\) were selected from absolute local calibration constants; their choices are independent of \(k\), the public exponents, the sign vector, the labels, the covariate, and the fair effect \(t\).
% lean: CausalSmith.Stat.LogoddsLowsmoothFrontier.SingletonNeighbourhood
% lean: CausalSmith.Stat.LogoddsLowsmoothFrontier.calibrated_support
\end{enumerate}
\end{proof}
